\chapter{Evaluation}\label{ch:eval}

% [SKELETON] Methodology first, then analysis of the results based on the methodology (merged from old 06+07).

The evaluation first discloses the methodology and, building on it, analyses the measurement series. Its
guiding question is cacheline awareness: Which design decision changes the cache-near measured
quantities --- L1, L2, and last-level misses, dTLB misses, branch misses, and the latency per
genus-interface call --- on which platform, and to what degree? Three evaluation objects structure the
chapter: the \emph{design space} (the hypotheses H1--H3 and the axis sensitivity analysis), the
\emph{apparatus} itself (the heuristic counter-proof and the measurement error that the measurement
instances introduce), and the \emph{automation} of the evaluation from the experiment XML, for at its
core the thesis automates precisely this evaluation. The methodology part (Section~\ref{sec:eval-method})
unfolds, in eight subsections, hypotheses and measured quantities, the XML-driven chain, workloads and
datasets, platforms and regimes, measurement levels and counters, build scope, fairness, and the
counter-proof; the results part (Section~\ref{sec:eval-results}) presents, in seven subsections, the
identity of the results, the data flow, the measurement series, the available measurements, the
sensitivity analysis, the validity check, and the summary.

The evaluation follows the chain customary in empirical software and systems research of research
question, hypothesis, measured quantity, experimental setup, execution, result, validity check, and
discussion~\cite{wohlin2012experimentation,shaw2003writing}: every hypothesis names its measured
quantity and its acceptance criterion before the measurement, every experimental setup its platform,
its regime, and its repetitions, and every result its
uncertainty~\cite{georges2007rigorous,kalibera2013rigorous}. Because the object is an artefact ---
the measurement apparatus with its axes ---, the assessment of the artefact according to the rules of
design-science research\footnote{Design science: the research direction of information systems that
makes a constructed artefact (here the measurement apparatus) itself the object of assessment and
demands for it utility, rigour of procedure, and traceability~\cite{hevner2004designscience}.}
is added: utility, rigour, and
traceability~\cite{hevner2004designscience,peffers2007dsrm,wieringa2014designscience}.
Performance measurements are reported according to the reporting rules of Hoefler and
Belli~\cite{hoefler2015benchmarking} and Blackburn et~al.~\cite{blackburn2016truth}: reference
quantity, repetitions, dispersion, platform, and values not collected (n/a) stand at every result. The
three evaluation objects structure the subsections; cacheline awareness is their common guiding axis.

The guiding axis of the evaluation is cacheline awareness: Compared are the cache-near measured
quantities --- L1, L2, and last-level misses, dTLB misses, branch misses, energy, and the latency
percentiles per genus-interface function; they are sixteen registered measurement categories with one
regime ordinal per category --- per axis assignment, per family- and genus-interface call (the channel
address always reads axis $\to$ genus $\to$ category), and per platform, that is, per registered machine
signature: prod1 (AMD Zen~5), prod2 (Intel Core i9-12900K, Alder Lake), and the ODROID board with Intel
Gracemont cores. The question is not which algorithm is faster, but which axis moves the
cacheline-related costs on which microarchitecture. The PMC categories (performance monitoring
counters, the hardware counters of the processors) are the primary measured quantities for this, the
time categories the cross-check; because hardware counters are trustworthy only after a cross-check by
measurement~\cite{weaver2008counters}, Section~\ref{sec:eval-counters} discloses the semantics of every
counter per model.

\section{Methodology}\label{sec:eval-method}
\subsection{Hypotheses and Measured Quantities}\label{sec:eval-hypotheses}
Three \emph{hypotheses} guide the analysis; each names its measured quantity and its acceptance
criterion before the measurement:
\begin{description}
  \item[H1 (page-type cost)] The mean access cost per page-type variant differs systematically by
  workload: dense page types win, as expected, under Zipf distribution (YCSB-A), range-optimized pages
  under range scans (YCSB-E). The measured quantities are the latency percentiles and the cache-miss
  categories per workload; H1 counts as supported if the ranking of the page-type variants flips
  between YCSB-A and YCSB-E on the same platform and the margin exceeds the uncertainty band of the
  build variants (Section~\ref{sec:eval-counterproof}).
  \item[H2 (code quality per source)] The code-quality score inherited from the original source
  repository correlates measurably with the achieved throughput. Two instruments supply the score:
  first, the documented manual audit per source repository over seven assessment axes (style, tests,
  documentation, maintainability, maintenance status, license, build system; scale 1--5) for 22
  assessed sources --- a record among the audits of the cache-engine repository ---, second, the
  tool-computed instrument that actually enters the measurement row: the weighted
  \texttt{cppcheck} finding density per thousand physical lines (weights error~3, warning~2,
  performance~1, portability~1, style~0.5) over the vendored original sources; the record lists 33
  entries, 22 of them as n/a because the reconstruction possesses no vendored original source, and the
  value travels as column \texttt{h2\_code\_quality\_score} of the WIDE schema (kept in the code as
  field \texttt{h2\_score}, record \texttt{sota\_h2\_scores.xml}) in every measurement row of a
  SOTA series. Both are complemented
  by the documented per-function \texttt{is\_original} marking (Appendix~\ref{app:blocks}).
  Acceptance criterion: H2 counts as supported if the tool-computed score over the substantiated SOTA
  profiles shows, with the throughput in the same cell (platform, workload, regime), a rank-correlation
  coefficient beyond a threshold named before the measurement; n/a profiles do not enter the
  correlation, and the manual audit is the qualitative classification.
  \item[H3 (inline vs.\ external vs.\ chain distribution)] The observed distribution of the three
  ValueHandle variants correlates with page density: dense pages prefer inline, sparse pages external,
  and PRT-ART heuristics ChainRef. The registry of the axis \texttt{value\_handle} lists five
  building blocks (Inline, External-Pool, Immutable-Shared-Ref, Versioned-Pointer, Chain-Ref; only
  Inline is bound to the golden reference catalogue); H3 regards the three variants decisive for
  PRT-ART as a reasoned excerpt. The measured quantity is the share per variant over page density; H3
  counts as supported if the share moves monotonically with density.
\end{description}
These \emph{evaluation hypotheses} are not to be confused with the identically named PRT-ART telemetry
metrics H1--H3 (Section~\ref{sec:prtart-demo}): the former are testable statements about the design
space, the latter measurement organs of the probe that, among other things, supply the data for this
test.

\subsection{Fully Automatic Experiment Evaluation from the XML}\label{sec:eval-xml}
The evaluation is itself an object of the automation (Figure~\ref{fig:eval-xml-kette}): A user XML
with the root \texttt{<comdare\_experiment>} declares \emph{what} is measured, \emph{where}, and
\emph{when} --- never the \emph{how}. Decisive for the evaluation are the sections
\texttt{workloads}, \texttt{datasets}, \texttt{measurement\_tooling}, \texttt{measurement\_categories},
\texttt{run\_methodology}, \texttt{writeback\_methods}, and \texttt{output}; the profile validator
checks them before every run against the registries (axis names and values, merge bindings, fairness
modes, records). From the XML the plan arises deterministically --- the director walk of the planner,
retrievable as a pure text emission even without a build ---, from it the build jobs of the
CacheEngineBuilder and the tier binaries, whose measurement is deposited in the workbook (xlsx); the
CSV form is the ordered tree of CSV files generated from the workbook (one folder per sheet, therein the
files of the individual sub-axis sweeps; not a single child document), and the
generators of the tool repository (binary-to-CSV, CSV-to-LaTeX, diagram generator, appendix
generator) render fact sheets, tables, and TikZ diagrams per language into a \texttt{tabellen}
directory that the manuscript includes. No step of this chain is executed by hand; every number in the
appendix carries the run that produced it (Appendix~\ref{app:measurements}).
Implementation status: The chain is demonstrated in a first end-to-end run --- the smoke measurement
series over 43 permutations in the appendix is its product ---; the XML section \texttt{<publish>} for
the store write-back, the command-line-guided creation mode of the XML
(Section~\ref{sec:impl-pipeline}), and the third output factory LaTeX
(Section~\ref{sec:eval-dataflow}) are planned but not implemented.
\begin{figure}[htbp]\centering
\begin{tikzpicture}[font=\footnotesize,>=Stealth,
  st/.style={draw,rounded corners,fill=blue!7,text width=2.55cm,minimum height=1.1cm,align=center,inner sep=2pt,font=\scriptsize},
  og/.style={st,fill=green!9},
  pl/.style={st,fill=orange!10,dashed,minimum height=0.8cm,font=\tiny}]
  \node[st] (x) at (0,0){\textbf{User XML}\\\texttt{<comdare\_\allowbreak{}experiment>}: workloads, datasets, tooling, categories, run\_methodology, writeback, output};
  \node[st] (v) at (3.1,0){\textbf{Validator}\\registries, records, merge binding, fairness modes};
  \node[st] (p) at (6.2,0){\textbf{Planner}\\director walk, plan text, build legend};
  \node[st] (c) at (9.3,0){\textbf{CEB}\\build jobs per system permutation, tier binaries};
  \node[st] (m) at (12.4,0){\textbf{Measurement}\\wallclock / macro / micro, PMC counters, one measurement thread};
  \node[og] (xl) at (12.4,-2.6){\textbf{Workbook (xlsx)}\\trunk: always arises; one data sheet (constant sheet key); sheet per sub-axis permutation: planned};
  \node[og] (cs) at (9.3,-2.6){\textbf{CSV tree} (tree: planned)\\folder per sheet; WIDE 189 columns, pipeline view 16 + 16};
  \node[og] (g) at (6.2,-2.6){\textbf{Generators}\\binary-to-CSV, CSV-to-LaTeX, diagrams, appendix};
  \node[og] (t) at (3.1,-2.6){\textbf{tabellen/}\\per language: fact sheets, booktabs tables, TikZ};
  \node[og] (ms) at (0,-2.6){\textbf{Manuscript}\\\textbackslash{}input in the appendix; every number with its run};
  \foreach \a/\b in {x/v,v/p,p/c,c/m}{\draw[->,thick] (\a)--(\b);}
  \draw[->,thick] (m)--(xl);
  \foreach \a/\b in {xl/cs,cs/g,g/t,t/ms}{\draw[->,thick] (\a)--(\b);}
  \node[pl] (a1) at (0,-4.5){not implemented: XML creation mode (command line, hardware filter)};
  \node[pl] (a2) at (12.4,-4.5){not implemented: \texttt{<publish>} section (store write-back)};
  \node[pl] (a3) at (9.3,-4.5){not implemented: third factory LaTeX (sheet order, measurement series per binary)};
  \draw[->,dashed] (a1)--(ms);
  \draw[->,dashed] (a2)--(xl);
  \draw[->,dashed] (a3)--(cs);
\end{tikzpicture}
\caption{The automated evaluation chain from the user XML to the manuscript. Upper row: The XML
declares what, where, and when; the validator checks it against the registries; the planner emits the
deterministic plan; the CEB builds the tier binaries per system permutation; the measurement collects
the three measurement levels and the counters with one measurement thread. Lower row: The workbook is
the trunk (one data sheet under a constant sheet key; the fan-out into one sheet per sub-axis
permutation is planned but not implemented), the CSV form its ordered tree of CSV files generated from
it (one folder per sheet; not a single child document; implemented is a flat file), the generators
render per language into the \texttt{tabellen} directory, the manuscript includes the files --- every
number carries its run. Dashed: parts not implemented (XML creation mode, \texttt{<publish>} section,
third factory).}\label{fig:eval-xml-kette}
\end{figure}

\subsection{Workloads and Datasets}\label{sec:eval-workloads}
Phase~5 of the measurement pipeline (Execute, Section~\ref{sec:impl-pipeline}) derives the YCSB
workload per permutation as a fallback heuristic from the profile's traversal tag
(Table~\ref{tab:workload-routing}), unless an \texttt{<expected\_workload>} tag of the profile
overrides the heuristic (Section~\ref{sec:sota-instances}); the routing matches by substring
(\texttt{RANGE}, \texttt{HOT\_PATH}) and not over a closed tag list, and the value of the
\texttt{<expected\_workload>} tag is resolved against the six tokens \texttt{YCSB\_A} to
\texttt{YCSB\_F}. The global default workload is the permutation over all available workloads; the
runtime dimension \texttt{workload} is registered for this as a dynamic dimension of the measurement
registry, fed from the load framework YCSB~\cite{cooper2010ycsb}. Implementation status: For
profile-less modules the driver falls back to exactly one workload, the run's workload option; the
full permutation is planned but not implemented. The data basis is real string datasets of differing
characteristics (Table~\ref{tab:datasets}), each with a reproducibility record of source, checksum,
line count, preprocessing, and seed rule for the synthetic miss, update, and delete keys; four of the
six datasets (\texttt{url}, \texttt{protein}, \texttt{tpcds-id}, \texttt{trec-terms}) possess a
record actually computed against the preprocessed files, \texttt{dna} the record of the
Pizza\&Chili source, and \texttt{xml} possesses a specification without checksum because the file is
not available locally --- there is no synthetic checksum (n/a). The records lie as XML files per
dataset in the tool repository, and a canon test checks the computed records against the six-element
canon of the table.

\begin{table}[!htbp]
  \centering
  \caption{Profile-aware workload routing (fallback heuristic)}\label{tab:workload-routing}
  \begin{tabular}{lll}
    \toprule
    \textbf{Traversal tag} & \textbf{YCSB workload} & \textbf{Characteristic} \\
    \midrule
    \texttt{RANGE\_SCAN} / \texttt{RANGE\_HOT\_PATH} & YCSB-E & range queries \\
    \texttt{HOT\_PATH} / \texttt{PRTART\_HOT\_PATH} & YCSB-A & 50/50 read/update, Zipf \\
    \texttt{STANDARD} (default) & YCSB-C & read-only \\
    \bottomrule
  \end{tabular}
\end{table}

\begin{table}[!htbp]
  \centering\small
  \caption{Real datasets, their characteristics, and their reproducibility record}\label{tab:datasets}
  \begin{tabular}{lll}
    \toprule
    \textbf{Dataset} & \textbf{Characteristic} & \textbf{Record} \\
    \midrule
    \texttt{url} & very long keys, long common prefixes, medium alphabet & real \\
    \texttt{dna} & small alphabet, high repetition, deep prefix entropy & Pizza\&Chili record \\
    \texttt{protein} & biological sequences, small to medium alphabet & real \\
    \texttt{xml} & structured strings with recurring prefixes & n/a (file not local) \\
    \texttt{tpcds-id} & database-like IDs, sorted and unsorted & real \\
    \texttt{trec-terms} & dictionary/search-index terms with defined miss keys & real \\
    \bottomrule
  \end{tabular}
\end{table}

\subsection{Platforms, Vendor Lanes, and Regimes}\label{sec:eval-platforms}
The measurements target the registered measurement machines of the apparatus: prod1, the AMD machine
with a Ryzen~9~9950X3D (Zen~5; machine signature \texttt{prod1\_zen5} with 22 features including AVX-512), % chktex 29
at the same time the development workstation of the verification; prod2, the Intel machine with a
Core i9-12900K (Alder Lake; signature \texttt{prod2\_alder\_lake} with nine features, without
AVX-512, with P and E cores); and an ODROID board with Intel Gracemont cores (signature
\texttt{odroid\_gracemont}, nine features). The ZIH cluster of TU Dresden is declared as sub-axis IS3
of the target ISA (Barnard, Capella, GraceHopper, N/A); the Barnard node possesses Intel Sapphire
Rapids CPUs, two Xeon Platinum 8470 per node~\cite{zih_hpc_compendium_hardware}, and Capella is the
GPU cluster with AMD Genoa CPUs and NVIDIA H100. In no machine signature is a ZIH node registered, and
a run there is not measured (n/a) --- in accordance with Section~\ref{sec:measurement-system}. An HBM
platform is registered in no signature and in no axis as a platform; the platform probe carries an HBM
feature and the allocator axis knows PIM special cases, nothing more --- it is planned but not
registered. An \texttt{IPlatformProbe} (Phase~1) reports per platform the runnable subset of the axes
(for instance AVX-512 only where available), whereby the exclusion of an unavailable extension is
always decided \emph{machine-side} --- with a warning and a named error class, never as a silent
profile filter, for the experiment profile stays identical for all machines; per platform a scalar and
an AVX2 build are prepared, an AVX-512 build only where the signature carries the feature. The
available axis options are checkable per platform via the planner probe; a command-line-guided XML
creation mode with platform filter is planned but not implemented
(Section~\ref{sec:impl-pipeline}). On hybrid CPUs, P and E cores are measured separately rather than
aggregated together, on Intel as on AMD\@; the measurement mechanics for this are the core PMU
domains of the kernel (on Intel the sysfs names \texttt{cpu\_core} and \texttt{cpu\_atom}), while the
naming of the sub-axis stays vendor-neutral --- the sysfs names stand explicitly on the prohibition
list of the ISA sub-axis labels. The core modes follow the six-mode model from
Section~\ref{sec:axis-triad}: four pinning modes per core type (P or E cores, a single core or all
cores of the type), the unpinned reference run, and \emph{HybridAware} as the sixth mode (scheduling
dimension \texttt{hetero\_core\_dispatch}), which is mandatory for the cache-line-awareness
measurement; the attachment of the core pinning and the core modes to the run profile is planned but
not implemented (Section~\ref{sec:impl-system-axes}).

The productive measurement operation is a \emph{two-machine operation} --- the AMD machine prod1
alongside the Intel machine prod2 ---, and the counter-based collection carries for this two vendor
lanes, each individually activatable, \texttt{pmc:amd} on prod1 and \texttt{pmc:intel} on prod2, each
with its own PMU resource group; the lanes deliberately run in parallel, because different hosts
possess different PMUs and a shared group would serialize them without reason. Completeness is an
evaluation rule, not a CI gate: A measurement series counts as complete only once both lanes have
proven and delivered their counter prerequisites; a lane run on one side only is reported as a partial
measurement, never silently carried as a whole (Section~\ref{sec:eval-validity}). The measurement plan
moreover provides, on the production target machines, \emph{two operating-system regimes} for every
measurement, to separate production fidelity from measurement depth: An \emph{immutable} operating
system (Talos) mirrors production use and enforces boot-/compile-static cache settings
(Section~\ref{sec:hardware-history}); a \emph{root Linux with full hardware-counter access}
(\texttt{perf} and the model-specific registers, MSR), by contrast, is what opens access in the first
place to the counter-based measurement categories (Section~\ref{sec:sota-instances}) and to the
runtime-configurable cache-near axis parameters (prefetch distance, hardware prefetcher). The second
regime is \emph{designed, not carried out}: The reported measurements all originate from the
root-Linux regime of the two production machines; the Talos regime is planned but not measured.

\subsection{Measurement Levels, Collectors, and Counters}\label{sec:eval-counters}
To be separated from this is which counters that access actually populates
(Table~\ref{tab:eval-pmc-zaehler}): The measurement chain opens via
\texttt{perf\_event\_open}~\cite{linux_perf_event_open,demelo2010perf} four generic counters ---
L1D misses, last-level misses, dTLB misses, and branch misses
(\texttt{PERF\_COUNT\_HW\_BRANCH\_MISSES}, portable on Intel and AMD). For L2 misses and coherence
invalidations there is no portable generic counter; both are collected model-bound via a RAW event
catalogue whose encodings are cross-checked: For
Zen~5\footnote{Zen~5: the AMD microarchitecture of the Ryzen~9~9950X3D (Granite Ridge, CPUID family 26); % chktex 29
prod1 possesses two L3 complexes of unequal size.} \texttt{l2\_cache\_req\_stat.ic\_dc\_miss\_in\_l2} is
the event \texttt{0x64} with mask \texttt{0x09} (demand misses in L2 without L2 prefetch) and % chktex 29
\texttt{ls\_dmnd\_fills\_from\_sys.remote\_cache} the event \texttt{0x43} with mask \texttt{0x14} % chktex 29
(demand fills from the cache of another CCX, a core complex with its own L3); a literal invalidation
counter does not exist in the core event set, the column carries exactly this assignment. The
encodings stem from the event table of the Linux kernel for Zen~5~\cite{linux_perf_amdzen5}, from
which \texttt{perf} binds the named events. Models without a catalogue entry carry n/a, never a
guessed encoding; for Intel no entry exists, because the cross-check on an Intel host is not
available. The last-level event fails on Zen~5 at opening (\texttt{ENOENT}). IPC/CPI is declared as a
category but not wired up --- the counter bundle lacks the sources for instructions and cycles ---,
and the energy draw is read best-effort from the RAPL sysfs zone\footnote{RAPL (Running Average Power
Limit): the energy counters of the processors, which the operating system reads out via the powercap
interface in microjoules~\cite{david2010rapl,khan2018rapl}; without read permission on the zone the
value stays empty.} and stays empty without read permission on the zone. The principle of explicitly
coded non-observation is realized as a field flag: Per counter the counter bundle carries a
\texttt{*\_source\_available} flag, seven flags for seven counters; a number means that the source
has delivered --- a measured zero included ---, n/a means that the source is missing, and this
distinction travels end-to-end into workbook, tables, and diagrams. Hardware counters thereby count
as trustworthy only after cross-checking~\cite{weaver2008counters}; the cross-check for Zen~5 is
substantiated by measurement, the one for Intel is not available.
\begin{table}[htbp]\centering\footnotesize
\begin{tabular}{@{}l>{\raggedright\arraybackslash}p{4.6cm}>{\raggedright\arraybackslash}p{3.2cm}>{\raggedright\arraybackslash}p{2.3cm}@{}}
\toprule
Category (column) & Source & Intel status & Validity flag \\
\midrule
\texttt{cache\_misses\_l1} & generic perf event (L1D read miss) & generic: yes & own flag \\
\texttt{cache\_misses\_l2} & RAW catalogue Zen~5: event \texttt{0x64}, mask \texttt{0x09} (demand misses in L2) & no catalogue entry, cross-check not performed (n/a) & own flag \\ % chktex 29
\texttt{cache\_misses\_l3} & generic perf event (last level); fails on Zen~5 at opening (n/a) & generic: yes & own flag \\
\texttt{dtlb\_misses} & generic perf event (dTLB read miss) & generic: yes & own flag \\
\texttt{branch\_misses} & generic perf event (\texttt{PERF\_COUNT\_\allowbreak{}HW\_\allowbreak{}BRANCH\_\allowbreak{}MISSES}) & generic: yes & own flag \\
\texttt{coherence\_invalidations} & RAW catalogue Zen~5: event \texttt{0x43}, mask \texttt{0x14} (foreign-CCX fills) & no catalogue entry, cross-check not performed (n/a) & own flag \\ % chktex 29
\texttt{energy\_micro\_joules} & RAPL sysfs zone, best-effort (read permission) & as AMD (read permission) & own flag \\
IPC/CPI & declared, not wired up (source missing) & n/a & --- \\
\bottomrule
\end{tabular}
\caption{The counters of the measurement chain and their sources: four generic perf events
(portable), two model-bound RAW events from the catalogue for Zen~5 (cross-checked by measurement),
the RAPL energy as best-effort reading, and IPC/CPI as a declared, not wired-up category. Per counter
the counter bundle carries a \texttt{*\_source\_available} flag; a number means \enquote{source has
delivered}, n/a means \enquote{source is missing}. The Intel status names what is collected on prod2
without a catalogue entry.}\label{tab:eval-pmc-zaehler}
\end{table}

Time and observer measurement are united in the measurement checkpoint
(\texttt{checkpoint\_measure}); their responsibilities are separated by the measurement-tooling main
axis in exactly this order: \emph{wallclock} (w) --- the wall clock of the overall runs in the
CacheEngineBuilder over the load frameworks, parameter-free; \emph{macro} (ma) --- the time of the
genus-interface functions of the tier binary; \emph{micro} (mi) --- the time and the success
parameters per axis interface (Section~\ref{sec:measurement-system},
Figure~\ref{fig:traeger-flaechen}); subsets of the three levels are allowed, every other order is
syntactically wrong, and w/ma/mi is the registered alias. Three collectors take over the collection
as building blocks of the measurement axis \texttt{collector} --- wall clock, observer snapshot, and
PMC ---, and the sixteen measurement categories are the sub-axis of the tooling axis, not its main
fan-out. The analysis keeps the three levels apart: axis interface, family and genus interface, and
the wall clock of the CacheEngineBuilder are never mixed (Section~\ref{sec:eval-sensitivity}); PMC
measurements as a figure type of their own in the appendix generator are planned but not implemented.
Implementation status: Registry and level canon are implemented; the order validation at the three
serialization points and the collector as a self-contained component are not.

The time- and observer-based categories run identically under both regimes and secure comparability;
the PMC categories are collected only under the privileged regime. Per category the registry carries
the regime as an ordinal (0 for time and observer, 1 for counters): nine time/observer categories and
seven counter categories. Operating system, kernel, and counter domain are to be logged per platform:
The operating-system probe collects kernel and identifier of the system as a system axis, and the
counter source records per counter on which core PMU domain it opened (generic or, on hybrid
machines, the sysfs name of the domain); a column of their own in the measurement output these three
items do not possess in the WIDE schema --- the columns \texttt{platform} and
\texttt{build\_version} carry platform tag and build version per row, the log column per domain is
not implemented.

\subsection{Build Scope and Optimization Levels}\label{sec:eval-build-scope}
The scope of the full build is fixed in the profile. The profile declares the build control
quantities: two optimization levels (O2, O3) times two SIMD levels (no extension, AVX2), that is,
four build variants per permutation (Section~\ref{sec:axis-triad}). Both levels are declared in the
system registry as IEEE-754-deterministic; the registry carries five levels (O0 to O3 and Ofast), and
excluded from the profile is Ofast alone, because it would break the floating-point determinism
according to IEEE~754 and the run-to-run determinism of the apparatus as well as the CRC-64 checksum
of the golden reference catalogue --- the registry marks exactly this level as non-deterministic.
Maximum optimization thus remains selectable: The default of the build is globally O2, O3 is selected
via the XML and built alongside, and the selected level is identity-effective via the toolchain
member of the fingerprint, so that O2 and O3 stocks are never tacitly confused as cache hits; O0 to
O3 as a user choice in the XML and the pass-through of further compiler flags are planned but not
implemented. Per build variant the golden reference catalogue is built: $2^{17} = 131072$ binary
identities, because the golden profile activates for seventeen organ main axes two representatives
each from the offering of the organ registry and the eighteenth (\texttt{persistence\_target}) is
fixed (Section~\ref{sec:mess-chain}); the system axes multiply the build matrix, never the
\texttt{binary\_id}. The cardinality of the real permutation space the planner computes dynamically
from its inputs (subcommand \texttt{simulate}); it is not a fixed number. Results of the runs of the
full build are not available (Section~\ref{sec:eval-available}).

\subsection{Fairness and Reproducibility}\label{sec:fairness}
Fair comparability is secured by fixed rules. Each comparison separates a common minimal mode (common
denominator: external value handles, no PRT-ART special paths) from the PRT-ART native mode (inline
switching and adaptive page-type selection per branch point); the user XML declares the mode per SOTA
series as \texttt{fairness\_mode} with the two permitted values \texttt{common\_denominator} and
\texttt{native}, and the validator rejects every other value. Implementation status: The mode is a
series metadatum without influence on the measurement --- it travels as a tag column of the
measurement row and in the resumption stamp, never in the \texttt{binary\_id} ---, the composition
definition of the minimal mode is data-bound, and two series of the same probe with different modes
therefore produce the same \texttt{binary\_id}, which the validator reports as a warning. All methods
receive identical, normalized byte keys as well as identical miss, update, and delete sets; the
warm-up is the two-phase driver of the load-profile orchestrator (save, warm up, roll back, measure),
otherwise measurement is cold, and each configuration is measured in several separate runs
(repetition dimension \texttt{repetition}). Thread pinning: The actuator --- a range-bound pinning of
the measurement thread --- is implemented; the core choice within the core complex is actuator policy
without an emission contract, and the measurement emission of the planner reports the pinning as
planned and the actual state as \emph{unpinned}, because the actuator is not attached to the run
profile --- on prod1 with its two L3 complexes of unequal size (96 and 32~MiB) an unpinned run is not
reproducible, which the emission says literally. Missing capabilities of a baseline --- for instance
range scans of a hash table --- are reported as n/a rather than covertly emulated; the mechanics for
this are the sample statuses \emph{NotApplicable}, \emph{SourceUnavailable}, and \emph{Failed} of
the measurement axes, so that missing capabilities and missing sources carry n/a, never a zero.
Provenance: toolchain (compiler identifier and version), flags, ISA summary, and the commit states of
the four repositories are captured at configuration time as build provenance and written into the
provenance manifest of the host exports; the allocator is, as an organ axis, part of every
\texttt{binary\_id}; a provenance per individual foreign library --- its own compiler, its flags, its
commit hash --- is not logged (not implemented). The latency percentiles are determined by the
evaluation tree of the cache engine nearest-rank from the raw values --- Hyndman and Fan, type~1:
rank $k = \lceil q \cdot n \rceil - 1$, without interpolation~\cite{hyndman1996quantiles} ---, so
that every percentile is an actually measured value\footnote{Nearest rank: The $q$-quantile is the
smallest sample value below which at least the share $q$ of the sample lies; interpolating methods
(type~7), by contrast, produce values that were never measured.}; the presentation tools of the tool
repository possess their own copies of the conversion (a declared gap), and the HDR
histogram~\cite{tene_hdrhistogram} is a method of its own of the distributional view
(Section~\ref{sec:impl-system-axes}), not a complement of the nearest rank. Layout and environment
effects that produce wrong conclusions without a recognizable error~\cite{mytkowicz2009wrongdata}
the apparatus counters with identical keys, identical build matrix per variant, and the repetition
dimension; a randomized layout stabilization in the sense of Curtsinger and
Berger~\cite{curtsinger2013stabilizer} is not part of the apparatus and is carried in
Section~\ref{sec:eval-validity} as a limit.

For foreign libraries the original-compiler rule applies: Dissected axis functions are bound in at
compile time and compiled into a once-again original binary block, so that every dissected
sub-algorithm keeps its performance properties; where feasible or comparable, the foreign source is
built with its original compiler. The proof of originality is implemented: the \texttt{is\_original}
marking per function (a macro of the cache engine together with the mixin base of the axis building
blocks), the validator as a program of its own, and the manual relock gate of the CI, which
re-anchors the lock files of the markings; the evidence per building block stands in
Appendix~\ref{app:blocks}. The H2 record marks reconstructions without a vendored original source as
n/a (22 of 33, Section~\ref{sec:eval-hypotheses}). Implementation status: The automatic intake of the
vendored sources into the build with their original compiler is planned but not implemented.

The measurement \emph{operation}, too, is part of these fairness rules, and it runs through the
store. The store is the stock of all built tier binaries and their measured values, addressed via
compile-time keys (Section~\ref{sec:lager}; Figure~\ref{fig:lager-bestaende} shows its structure):
stock~1 keeps the binary under its SHA-512 fingerprint, stock~2 the measured values under the pair of
the fingerprint of the measured binary and the hardware identity of the measuring machine; the cell
of ISA, optimization, and extension stands bracketed beside it, never in the digest. Because the
measurement-gate state is a fingerprint member, the binary with and the binary without measurement
sensors lie side by side as two entries of stock~1. Measurement is always switched on; in the store
only what is missing there is re-measured: A binary found is skipped if its fingerprint is known and
its passed test attestation is inventoried at the binary, otherwise it is built or measured
(Figure~\ref{fig:eval-lager-kreislauf}). The CEB controls the parallelism by producing parallel
requests to the tier binaries --- the execution threads arise in the CEB, in a worker pool with
deterministic result order ---, while the threading of the family stays in the tier binary.
Implementation status: Stock~1 and~2 as well as the key schema of all four stocks are implemented;
stock~3 and~4 (synthesized curves as files, solutions per user XML), the planner sheet function, the
full placement of the CEB, the hybrid storage path, the commit skip, the rebuild flag, and the XML
section \texttt{<publish>} form the full store build-out and are not implemented.
\begin{figure}[htbp]\centering
\begin{tikzpicture}[font=\footnotesize,>=Stealth,thick,
  k/.style={draw,rounded corners,fill=blue!7,text width=2.7cm,minimum height=1.15cm,align=center,inner sep=2pt,font=\scriptsize},
  lg/.style={k,fill=green!10},
  dc/.style={draw,fill=orange!14,text width=2.7cm,minimum height=1.15cm,align=center,inner sep=2pt,font=\scriptsize},
  pl/.style={k,fill=orange!10,dashed,font=\tiny,minimum height=0.9cm}]
  \node[k]  (bau)  at (0,0){\textbf{Build (CEB)}\\batches of 4096 binaries each, inventoried resumption};
  \node[lg] (b1)   at (4.1,0){\textbf{Stock 1: binary}\\key: SHA-512 fingerprint; cell bracketed beside it};
  \node[dc] (skip) at (8.2,0){\textbf{Recognition criterion}\\fingerprint known \emph{and} test attestation inventoried as passed?};
  \node[k]  (mess) at (12.3,0){\textbf{Measurement}\\probe dock; one measurement thread; the CEB produces parallel requests};
  \node[lg] (b2)   at (12.3,-2.6){\textbf{Stock 2: measured values}\\key: fingerprint $+$ hardware identity; cell beside it};
  \node[lg] (mp)   at (8.2,-2.6){\textbf{Workbook (xlsx)}\\CSV tree (folder per sheet; tree: planned); appendix generators};
  \node[k]  (ref)  at (4.1,-2.6){\textbf{Structure of the store}\\Figure~\ref{fig:lager-bestaende}: tree, stocks 1--4, key schema};
  \node[pl] (vp)   at (0,-2.6){full store build-out: stock 3/4, hybrid storage path, commit skip, rebuild flag, \texttt{<publish>}};
  \draw[->] (bau)--(b1);
  \draw[->] (b1)--(skip);
  \draw[->] (skip)--node[above,font=\tiny]{no: measure}(mess);
  \draw[->] (mess)--(b2);
  \draw[->] (b2)--(mp);
  \draw[->] (skip)--node[right,font=\tiny,align=left]{yes: skip,\\measured values exist}(mp);
  \draw[->] (skip.north) -- ++(0,0.75) -| node[pos=0.25,above,font=\tiny]{fingerprint unknown: build} (bau.north);
  \draw[->,dashed] (vp)--(ref);
\end{tikzpicture}
\caption{The store cycle in measurement operation: The CEB builds in batches of 4096 binaries each;
every binary lies under its SHA-512 fingerprint in stock~1; the recognition criterion skips what is
inventoried with a passed test attestation and otherwise lets build or measure; the measurement runs
at the probe dock with one measurement thread, while the CEB produces the parallel requests; the
measured values lie under fingerprint and hardware identity in stock~2 and flow into the workbook
(the CSV tree form with one folder per sheet is planned; implemented is a flat file). The structure
of the store (tree, four stocks, key schema) is shown by Figure~\ref{fig:lager-bestaende}; dashed the
full store build-out.}\label{fig:eval-lager-kreislauf}
\end{figure}

Measurement remains strictly separated from building: measurement runs never start automatically with
an error-free build but are triggered explicitly --- an error-free build proves buildability, not a
measurement result. The run methodology knows four work modes --- \texttt{build}, \texttt{measure},
\texttt{compare}, and \texttt{release} (Section~\ref{sec:impl-pipeline}) ---; the measure mode
measures with exactly \emph{one} measurement thread, a parallel measurement execution is armed via no
registry row and reachable only via the state injection of the run entry, and \texttt{--debug} is a
switch of the planner shell for the wiring check, which stays locked for users, not a work mode; the
threading of the family itself is thereby left to the tier binary rather than overridden by the
harness. A binary found in the store is skipped only if its passed test attestation is inventoried at
the binary (Figure~\ref{fig:eval-lager-kreislauf}), and every measurement is preceded by the
conformity check at the probe dock (Section~\ref{sec:impl-contracts}): the conformity gate of the
genus contract before the post-compile attestations of the measurement interfaces, beside it the
consistency gate of the measurement; build as well as measurement runs are processed per machine in
milestone batches of 4096 binaries each at a stretch --- the grain is a compile-time constant whose
deviation is a compile error ---, with an evenly distributed allocation of the jointly buildable
binaries across the machines and inventoried resumption after interruption.

\subsection{The Heuristic Counter-Proof as Fourth Dimension}\label{sec:eval-counterproof}
Beyond the three hypotheses, the analysis receives a fourth, self-contained \emph{evaluation
dimension}: the heuristic counter-proof. Its concept is developed in Chapter~\ref{ch:measurement}
(Sections~\ref{sec:heuristics-modes} and~\ref{sec:heuristics-curves}): The CacheEngineBuilder stage
compiles the heuristically estimated recombination of the tier binaries into a hybrid tier binary ---
the hybrid module emitter lies in the builder tree of the CEB ---, and the hybrid tier binary is a
carrier stage of its own \emph{behind} the CEB, which docks at the probe dock and itself produces no
code. The measurement-curve type system supplies the estimation basis
(Figure~\ref{fig:eval-messkurven}): per axis one curve over its parameters and per measurement level
wallclock, macro, and micro, from these the break-even points and the estimated recombination; the
measurement-error detection is formulated by Section~\ref{sec:heuristics-curves}. The assessment
question of this dimension reads: \enquote{Is the heuristically estimated recombination of the tier
binaries faster than the existing paper algorithms?} Unlike H1--H3, which test statements about the
design space, the counter-proof tests the apparatus of the thesis itself: whether the recombination
estimated from the measurement curves reaches, in actual wall-clock behaviour, the performance
predicted by the curves; it therefore stands \emph{beside} the ranking of the sensitivity analysis
(Section~\ref{sec:eval-sensitivity}), not within it. Implementation status: Classification layer,
dock layer, XML parser, and binary proxy of the hybrid stage are implemented (six of the nine files of
the dock, proxy, and router layer envisaged in the \texttt{README} of the directory
\texttt{hybrid/}, Section~\ref{sec:impl-hybrid-lager}); router and eviction presuppose measurement
curves and are not implemented.
\begin{figure}[htbp]\centering
\begin{tikzpicture}[font=\footnotesize,>=Stealth,
  k/.style={draw,rounded corners,fill=blue!7,text width=2.55cm,minimum height=1.1cm,align=center,inner sep=2pt,font=\scriptsize},
  hy/.style={k,fill=orange!14},
  pl/.style={k,fill=orange!10,dashed,font=\tiny,minimum height=0.8cm},
  bv/.style={draw,rounded corners,fill=green!9,text width=4.5cm,minimum height=1.0cm,align=center,inner sep=2pt,font=\tiny}]
  \draw[->] (0,0) -- (4.8,0) node[below left,font=\scriptsize]{axis parameter (e.g.\ working set)};
  \draw[->] (0,0) -- (0,3.2) node[above,font=\scriptsize]{metric per level};
  \draw[blue!70!black,thick] plot[smooth] coordinates {(0.3,0.5) (1.3,0.8) (2.3,1.5) (3.3,2.3) (4.3,2.9)};
  \draw[green!50!black,thick] plot[smooth] coordinates {(0.3,1.6) (1.3,1.5) (2.3,1.5) (3.3,1.6) (4.3,1.8)};
  \draw[blue!70!black,thin,dashed] plot[smooth] coordinates {(0.3,0.7) (1.3,1.0) (2.3,1.7) (3.3,2.5) (4.3,3.1)};
  \draw[green!50!black,thin,dashed] plot[smooth] coordinates {(0.3,1.8) (1.3,1.7) (2.3,1.7) (3.3,1.8) (4.3,2.0)};
  \fill[red!70!black] (2.3,1.5) circle (2pt);
  \node[font=\tiny,text=red!70!black,anchor=south east] at (2.25,1.6){break-even};
  \node[font=\tiny,text=blue!70!black,anchor=west] at (3.35,2.9){variant A};
  \node[font=\tiny,text=green!50!black,anchor=west] at (3.55,1.35){variant B};
  \node[font=\tiny,align=left,anchor=north west] at (0.15,3.15){one curve per axis;\\bands wallclock / macro / micro};
  \node[k]  (re) at (6.7,1.6){\textbf{Estimated recombination}\\break-even per axis and level; curves from the store};
  \node[k]  (em) at (9.7,1.6){\textbf{CEB}\\hybrid module emitter in the builder tree: compiles};
  \node[hy] (hb) at (12.7,1.6){\textbf{Hybrid tier binary}\\stage behind the CEB; docks at the probe dock; produces no code};
  \node[k]  (wc) at (12.7,-0.6){\textbf{Wall-clock comparison}\\primarily against paper reconstructions, secondarily against the best single binary};
  \node[pl] (ro) at (9.7,-0.6){router and eviction: not implemented (presuppose measurement curves)};
  \draw[->,thick] (4.8,1.6)--(re);
  \draw[->,thick] (re)--(em);
  \draw[->,thick] (em)--(hb);
  \draw[->,thick] (hb)--(wc);
  \draw[->,dashed] (ro)--(hb);
  \node[bv] (v1) at (2.35,-2.6){\textbf{Variant 1}: observers in all tier binaries and in the hybrid tier\\gate fields: tier \texttt{tm}/\texttt{tmi} on, hybrid \texttt{hm}/\texttt{hmi} on};
  \node[bv] (v2) at (7.35,-2.6){\textbf{Variant 2}: observers only in the hybrid tier\\gate fields: tier off, hybrid on};
  \node[bv] (v3) at (12.35,-2.6){\textbf{Variant 3}: observers in no compiled artefact, wall clock only\\gate fields: all off};
  \node[font=\tiny,align=center] at (7.35,-3.55){uncertainty band from the differences 1$\to$2 (measurement organs of the animals) and 2$\to$3 (measurement commands of the hybrid tier); validity criterion $t_1 \ge t_2 \ge t_3$};
\end{tikzpicture}
\caption{The measurement-curve type system in the heuristic counter-proof. Left: per axis one curve
of the metric over the axis parameter, separated by measurement level (wallclock, macro, micro); the
intersection of two variants is the break-even point from which the recombination is estimated.
Middle and right: The CEB compiles the estimated recombination via the hybrid module emitter; the
hybrid tier binary is the carrier stage behind the CEB, docks at the probe dock, and itself produces
no code; the comparison is made on the wall clock, primarily against the paper reconstructions,
secondarily against the best single binary. Bottom: the three build variants of the
measurement-error test plan, kept apart via the gate fields of the fingerprint. Dashed: router and
eviction (not implemented).}\label{fig:eval-messkurven}
\end{figure}

The test plan for this is fixed in two stages; the corresponding measurement runs are not available
(Section~\ref{sec:eval-available}). The first stage quantifies the measurement error that the
measurement instances themselves introduce, for the heuristic's estimation basis is \emph{observed}
curves, while its claim concerns unobserved production behaviour. To this end, the same experiment is
set up in three build variants:
\begin{enumerate}
  \item observers in \emph{all} tier binaries \emph{and} in the heuristic hybrid tier,
  \item observers \emph{only} in the heuristic tier binary,
  \item observers in \emph{none} of the compiled artefacts --- only the wall clock is measured.
\end{enumerate}
The variants are kept apart via the measurement-gate member of the fingerprint
(Section~\ref{sec:stamp-model}): It carries fields of its own for the tier level (\texttt{tm},
\texttt{tmi}) and the hybrid level (\texttt{hm}, \texttt{hmi}); the off state of all fields reads
\texttt{mg=m0;s0;st0;x0;tw0;tm0;tmi0;hm0;hmi0}, and because the member enters the digest, the three
variants lie side by side in the store as three binaries with three fingerprints. The differences
between the three variants quantify the observer overhead separated by its origin --- measurement
organs of the subordinate animals (difference from variant~1 to~2) versus measurement commands of the
hybrid tier (difference from variant~2 to~3) --- and yield, per workload and platform, an
\emph{uncertainty band} below which two configurations count as indistinguishable; the band is formed
from repetitions with a confidence interval, not from a single run~\cite{georges2007rigorous}.
Validity criterion of the stage: With measurement sensors it must be slower than without --- for the
wall-clock times of the three variants $t_1 \ge t_2 \ge t_3$ holds; a violation is a measurement
defect of the affected cell, not a result, and leads to the exclusion of the cell. Implementation
status: This plausibility check is not implemented in the code; it connects to the measurement-error
interpolation as the fifth evaluation component of the compare stage
(Section~\ref{sec:heuristics-curves}), whose constellations the planner computes dynamically from the
gate assignment, never as a fixed number.

The second stage is the complementary \emph{fourth step} of the final heuristic assembly: The tier
binaries subordinate to the heuristic tier --- not the finally optimized main binary --- are rebuilt
without micro- and macro-benchmarks and without observers, so that no measurement artefact remains in
the compiled artefact. The comparison is then made purely on the wall clock, in two comparison sets
with a fixed precedence: primarily against the reconstructed paper algorithms --- the 33 corpus
papers with one profile XML each, 30 full profiles and 3 abstract ones
(Section~\ref{sec:sota-instances}); the paper comparison takes precedence, for it answers the
assessment question literally ---, secondarily against all known tier-binary permutations, in
particular the best single binary; this second comparison classifies whether a potential lead is
attributable to the heuristic as a whole or merely to a single superior permutation.

From this follows the decision rule: The counter-proof counts as \emph{delivered} per workload class
if the observer-free rebuilt recombination beats the fastest paper reconstruction on the wall clock
and the margin exceeds the uncertainty band of the build variants; it counts as \emph{missed} if a
paper algorithm or the best single binary comes out ahead --- this outcome, too, would be a result,
for it would quantify for the first time how much estimation error lies between curve prognosis and
real behaviour. The uncertainty band is to be distinguished from a measurement defect: Cells that
violate the validity criterion of the first stage do not enter the decision. The measurement data of
this dimension are not available; they are reported with series A--C in
Section~\ref{sec:eval-available}.

\section{Results and Analysis}\label{sec:eval-results}
\subsection{Identity of the Results: Stamp, Store, Workbook}\label{sec:eval-identity}
Every result row is identifiable via the stamp system (Section~\ref{sec:stamp-model}): the
\texttt{binary\_id} --- the path of eighteen \texttt{axis=value} segments in slot order ---, the
SHA-512 fingerprint over eleven members with the format identifier \texttt{fingerprint\_format=6},
and the three compiled-in stamp lines per realm (organ, system, measurement). This identity travels
unchanged through the chain: as store key (stock~1 and~2), as columns \texttt{permutation\_id} and
\texttt{fingerprint} of the measurement row, as sheet key of the workbook, and as identifier in every
diagram and every table of the appendix --- thus every presented number remains traceable to its
binary, its axis versions, and its machine. The mixing strategy of a probe is expressed not via a
version suffix but via the \texttt{merge} tokens \texttt{replace}, \texttt{merge}, and
\texttt{union} of the user XML and the hierarchically extended algorithm names of the organ line
under the namespace prefix of the probe; the flag \texttt{e} of the entry grammar denotes an
efficiency-core algorithm (the E-core path of a hybrid CPU), not an experiment and not a mixture.
Values of the measurement never appear in the stamp; what is stamped is the compile-time identity.
Implementation status: The appendix tables carry \texttt{permutation\_id} and
\texttt{fingerprint} per row; the full stamp identity per diagram is planned but not implemented.

\subsection{Data Flow from the Measurement to the Manuscript}\label{sec:eval-dataflow}
Per series run the measurement stage writes the one official \texttt{measurements.csv} as a
projection of the workbook (implemented is a flat file; planned is the CSV tree with one folder per
sheet, below); the \texttt{<output>} paths of the user XML (pattern
\texttt{\_runs/<date>/<series>/}) the run entry resolves as run provenance. Three column lists are
to be distinguished (Table~\ref{tab:eval-schema-ebenen}): the WIDE schema with 189 columns (separator
semicolon; the schema authority of the measurement output, whose header arises at runtime from four
individual sources and is therefore anchored as a schema freeze), the pipeline view with 16 columns
plus 16 additional columns, that is 32, for the LaTeX stages (separator comma), and the
thirteen-column legacy form of the result aggregator over \texttt{ExecutionResult}. The pipeline
view carries per permutation \texttt{permutation\_id}, \texttt{fingerprint}, \texttt{succeeded},
\texttt{workload\_used}, \texttt{op\_count}, \texttt{total\_cycles}, \texttt{cache\_misses\_l1},
\texttt{cache\_misses\_l2}, \texttt{cache\_misses\_l3}, \texttt{dtlb\_misses},
\texttt{coherence\_\allowbreak{}invalidations},
\texttt{energy\_\allowbreak{}micro\_\allowbreak{}joules}, \texttt{bytes\_\allowbreak{}allocated}, \texttt{bytes\_\allowbreak{}in\_\allowbreak{}use\_\allowbreak{}peak},
\texttt{external\_\allowbreak{}frag}, and \texttt{internal\_\allowbreak{}frag}; the full view adds six search counters
(\texttt{search\_\allowbreak{}insert} to \texttt{search\_\allowbreak{}peak\_\allowbreak{}occupancy}), \texttt{pmc\_\allowbreak{}available},
\texttt{branch\_\allowbreak{}misses}, \texttt{throughput\_\allowbreak{}ops\_\allowbreak{}per\_\allowbreak{}sec}, and the seven
\texttt{*\_\allowbreak{}source\_\allowbreak{}available} flags. \texttt{total\_cycles} carries --- under its legacy name ---
the representative p50 latency in nanoseconds, not cycles; the renaming or an actual cycle channel is
a documented design decision. Values not collected stand as n/a, never as zero.
\begin{table}[htbp]\centering\footnotesize
\begin{tabular}{@{}l>{\raggedright\arraybackslash}p{1.5cm}>{\raggedright\arraybackslash}p{1.9cm}>{\raggedright\arraybackslash}p{4.6cm}>{\raggedright\arraybackslash}p{4.0cm}@{}}
\toprule
Schema & Separator & Columns & Authority / consumer & Source in the code \\
\midrule
WIDE (E4) & semicolon & 189 & measurement output; header at runtime from four individual sources, anchored as a schema freeze & iterator of the CacheEngineBuilder (\texttt{lazy\_csv\_header}); schema freeze \\
PIPELINE & comma & 16 (+16 = 32) & LaTeX stages 04/05/06; full view as prefix superset & pipeline CSV schema: one list, four writers \\
F15-LEGACY & comma & 13 (12 + 1) & result aggregator over \texttt{ExecutionResult} & result aggregator of the builder commands \\
\bottomrule
\end{tabular}
\caption{The three measurement column lists of the cache engine and their authority.
\texttt{total\_cycles} of the pipeline view carries the representative p50 latency in nanoseconds
(legacy name), not cycles. F15-LEGACY comprises twelve columns and, appended at the end,
\texttt{degeneriert} (12~+~1).}\label{tab:eval-schema-ebenen}
\end{table}

The workbook (xlsx) always arises --- in memory, unconditionally, even if a profile persists CSV
only; it is the output and the standard, and the CSV form is not a single child document but an
ordered tree of CSV files --- one overall folder per workbook, beneath it one folder per sheet,
therein the CSV files of the individual sub-axis sweeps ---, which is always generated from the
workbook (xlsx to csv, as a strategy of the same factory) and is, as a presentation, its analogously
mirrored counterpart: A row reaches the CSV tree only if the workbook has accepted it before, and if
the workbook generation breaks, the CSV tree is invalid too; what goes to disk is decided by
\texttt{<writeback\_methods>} alone --- xlsx, csv, or both (the workbook, the CSV tree, or both at
the same target), as a strategy over the output formats with xlsx as default
(Section~\ref{sec:impl-system-axes}); the xlsx writer lies in the store deposit and writes via
libxlsxwriter. Implementation status: The seam writes the CSV projection as one flat file
(\texttt{projektion\_\allowbreak{}oeffnen} in
\texttt{ergebnis\_\allowbreak{}mappe\_\allowbreak{}naht}); the tree form is planned but not
implemented. The analysis enriches the projection via binary-to-CSV with profile ID, paper reference,
and axis assignment, renders CSV-to-LaTeX fact sheets per permutation ID (an empty input produces no
file, never an empty plot) and diagram-generator diagrams as TikZ bars, scatter plots, and heatmaps
together with the families of the measurement appendix built on them (ECDF, range bars, sweep
curves, 3D raw-data view, ratio matrices, normalized bars, Pareto scatter, per-axis attribution,
observer detail table) with A4 cap and size check (a 1$\times$1, 1$\times$n, or n$\times$1 matrix
produces no image); the appendix generator writes per language into a \texttt{tabellen} directory,
from which the manuscript includes the files --- the German version includes 38 generated files of
the appendix.

The workbook of a run is one file with one sheet per selected sub-axis permutation and an info sheet
(system information of the measuring machine, main axes used of all three realms, constant meta
axes); the sheets are numbered in the binding order measurement $\to$ system $\to$ organ
(\texttt{S001} to \texttt{Snnn}), and the columns stem exclusively from the WIDE header; columns
whose value never changes within a run move into the metadata of the info sheet instead of being
repeated in every row. The contract tests of the persistence check at the CSV output; a check at the
standard xlsx is planned but not implemented (Section~\ref{sec:impl-contracts}). Besides xlsx and
csv, a third factory LaTeX is planned, which serializes the sheet order of the workbook when rolling
out the appendix and carries per binary one table row as a temporal measurement series.
Implementation status: Writer and projection are implemented; the seam writes one data sheet under a
constant sheet key, and the projection is a flat file. Not implemented are: the fan-out into one
sheet per sub-axis permutation plus info sheet (Section~\ref{sec:limitations}), the CSV tree (one
folder per sheet), the check at the xlsx produced per binary (the persistence check reads the CSV
output), the export element with target filter and per-binary xlsx, the chart functions, the LaTeX
factory, the warnings section of the info sheet, and the catch-up of the WIDE freeze.

\subsection{Measurement Series and Merge Strategies}\label{sec:eval-series}
Three measurement series are evaluated (Table~\ref{tab:stage-series}): series~A --- PRT-ART against
the state of the art with the evaluation hypotheses H1--H3 (Section~\ref{sec:eval-hypotheses});
series~B --- the systematic variation of the cache-engine permutations, fed from the union of
stage~3; series~C --- merge/regression of old against new, across builds. Orthogonal to this
assignment, the user XML declares per axis the \texttt{merge} mode of the probe
(Table~\ref{tab:eval-verbund-modi}): \texttt{replace} (the default, also with an empty attribute; the
probe axis replaces the cache-engine axis with fallback, merge strategy
\texttt{Verbund2\_Replace}), \texttt{merge} (cache engine and probe hybrid per probe,
\texttt{Verbund2\_Hybrid}), and \texttt{union} (the explicit union of all axis variants without
discarding, bound to stage~3, \texttt{Verbund3\_Union}); stage~1 is \texttt{Verbund1\_CeOnly}, the
absence of any probe directive, and \enquote{stage~1 to~3} remains the alias of the merge-strategy
names, the full join the alias of \texttt{union}. Series~A works with \texttt{replace} and
\texttt{merge}, series~B with \texttt{union}. The states remain distinguishable at the binary via
the hierarchically extended algorithm names of the organ line (Section~\ref{sec:eval-identity}); a
stamp line of their own they do not possess, for the probe mixes against the organ axes and thus
already stands in their line, and a version suffix does not denote the mixture. For pure replacement
builds as for pure cache-engine builds, the name extension consequently stays absent. Implementation
status: The token \texttt{merge} is registered in the merge plan and declared in the golden core
profile; materialized are the strategies \texttt{Verbund1\_CeOnly}, \texttt{Verbund2\_Replace}, and
\texttt{Verbund3\_Union}, while \texttt{Verbund2\_Hybrid} remains without a production caller,
because its materialization in the strategy enum is not implemented.
\begin{table}[htbp]\centering\footnotesize
\begin{tabular}{@{}ll>{\raggedright\arraybackslash}p{2.2cm}>{\raggedright\arraybackslash}p{4.4cm}>{\raggedright\arraybackslash}p{1.6cm}@{}}
\toprule
\texttt{merge} token & Merge strategy & Alias & Effect & Series \\
\midrule
--- (no directive) & \texttt{Verbund1\_CeOnly} & stage~1 & cache engine only: standard building blocks and SOTA profiles & A (baseline, also C) \\
empty / \texttt{replace} & \texttt{Verbund2\_Replace} & stage~2 (replace) & the probe axis replaces the cache-engine axis, with fallback & A \\
\texttt{merge} & \texttt{Verbund2\_Hybrid} & stage~2 (hybrid) & cache engine and probe hybrid per probe; materialization open, no production caller & A \\
\texttt{union} & \texttt{Verbund3\_Union} & stage~3 (full join) & union of all axis variants without discarding, bound to stage~3 & B \\
\bottomrule
\end{tabular}
\caption{The \texttt{merge} tokens of the user XML and their merge strategies in the merge plan; the
stage names are aliases, the full join is the alias of \texttt{union}. The series column follows the
assignment from Table~\ref{tab:stage-series}.}\label{tab:eval-verbund-modi}
\end{table}

\subsection{Available Measurements}\label{sec:eval-available}
Methodology and evaluation pipeline are in place and verified on real data: A smoke measurement
series over 43 permutations --- DLL measurement; 27 Cartesian combinations of SIMD, memory layout,
and allocator plus 16 more over four two-valued quantities --- is documented by
Appendix~\ref{app:measurements} together with its limitations
(Section~\ref{sec:measurements:limitations}); it varies seven quantities (SIMD is a system axis,
memory layout and allocator are organ axes, the remaining four names are not registry building
blocks and are not equated), the full measurement series varies all eighteen composition axes. The
complete runs of the mandatory measurement series A--C on the target platforms are not available
(n/a), because the full build is not carried out (Section~\ref{sec:limitations}); the same holds for
the measurement data of the counter-proof (Section~\ref{sec:eval-counterproof}). No number of this
chapter stems from any run other than the one reported in the appendix.

% Insertion after the first complete run of series A--C (generator output per language):
%   \input{anhang/en/tabellen/<name>.tex} -- series A_full (substantiated: the 33 experiment XMLs declare A_full);
%   the file names of series B and C arise with the run.

\subsection{Axis Sensitivity Analysis}\label{sec:eval-sensitivity}
The core of the evaluation is the \emph{axis sensitivity analysis}: It answers the separability
problem of the task (Sections~\ref{sec:rqs} and~\ref{sec:task-concerns}) in that the axis framework
isolates each constituent as an interchangeable axis over an otherwise constant structure and the
contribution of each design decision becomes attributable under all-else-equal conditions. The
ranking runs over the three measurement levels in canonical order: \emph{wallclock} --- the overall
benchmarking of the YCSB load profiles in the CacheEngineBuilder ---, \emph{macro} --- the
genus-interface operations of the tier binary --- and \emph{micro} --- the axis interface, the set of
functions that all algorithms of an axis must provide (Section~\ref{sec:mess-chain}); from the rows
per permutation the analysis quantifies at the micro level which axis contributes most to the
performance variance, the wallclock level checks the internal micro measurement from outside, and
the measurement categories are thereby the sub-axis of the tooling axis
(Section~\ref{sec:eval-counters}). The guiding axis is cacheline awareness: which axis moves the
cache-near categories on which platform. From the ranking follows the recommendation of standard
configurations per workload class; where ratios are aggregated over workloads or platforms, this is
done by geometric, never arithmetic, mean~\cite{fleming1986statistics}. The heuristic counter-proof
stands beside this ranking (Section~\ref{sec:eval-counterproof}).

\subsection{Validity and Limits}\label{sec:eval-validity}
The validity of the evaluation is checked according to the four classes of empirical software
research\footnote{Threats to validity: the classification, customary in empirical software research,
of the threats to a conclusion into construct validity (does the measured quantity measure what is
meant?), internal validity (is the effect causally isolated?), external validity (does the finding
transfer to other platforms and loads?), and conclusion validity (do statistics and aggregation
support the conclusion?)~\cite{wohlin2012experimentation}.} (Table~\ref{tab:eval-gueltigkeit};
Section~\ref{sec:measurement-system} carries the same classes as quality criteria of the measurement
system). Construct validity: Do the categories measure cacheline awareness? The PMC semantics are
disclosed per model (Section~\ref{sec:eval-counters}), and hardware counters count as trustworthy
only when cross-checked~\cite{weaver2008counters}. Internal validity: The observer overhead is
bounded via the uncertainty band of the three build variants (Section~\ref{sec:eval-counterproof});
layout and environment bias~\cite{mytkowicz2009wrongdata,curtsinger2013stabilizer} are countered by
identical keys, identical build matrix, and repetitions, a layout randomization is missing; warm-up
and repetitions follow Section~\ref{sec:fairness}~\cite{kalibera2013rigorous}. External validity: The
registered measurement machines comprise two x86-64 microarchitectures and one Gracemont node; ZIH
and HBM are declared, not measured (Section~\ref{sec:eval-platforms}), the Intel RAW assignment is
missing, and a lane run on one side only counts as a partial measurement. Conclusion validity:
percentiles nearest-rank without interpolation, geometric aggregation, n/a instead of zero. The
limitations of the measurement appendix (Table~\ref{tab:le:limitierung},
Section~\ref{sec:measurements:limitations}) and of the conclusion (Section~\ref{sec:limitations}) are
the substantiated list of these limits; no number is named without a run, and the reporting rules for
performance measurements~\cite{hoefler2015benchmarking,blackburn2016truth} apply to every number that
a run delivers.
\begin{table}[htbp]\centering\footnotesize
\begin{tabular}{@{}>{\raggedright\arraybackslash}p{1.6cm}>{\raggedright\arraybackslash}p{4.0cm}>{\raggedright\arraybackslash}p{4.9cm}>{\raggedright\arraybackslash}p{3.7cm}@{}}
\toprule
Class & Threat & Countermeasure in the apparatus & Implementation status \\
\midrule
Construct & counters do not measure what is meant (model semantics, prefetch shares) & RAW catalogue per model with cross-check; semantics disclosure per column; n/a without catalogue entry & implemented (Zen~5); Intel entry: not implemented \\
Construct & time category misread as cycles & \texttt{total\_cycles} declared as p50 latency in nanoseconds & implemented; renaming: design decision \\
Internal & observer overhead distorts the curves & three build variants, uncertainty band, validity criterion $t_1 \ge t_2 \ge t_3$ & gate fields implemented; check: not implemented \\
Internal & layout/environment bias, change of the L3 complex & identical keys and build matrix; repetition dimension; pinning reported as planned & actuator implemented, attachment: not implemented; no layout randomization \\
Internal & warm-up and cold-start effects & two-phase driver (save, warm up, roll back, measure) & implemented \\
External & measurement machines do not cover the target platforms & registered signatures prod1, prod2, ODROID\@; ZIH declared as IS3 & ZIH and HBM\@: declared, not measured (n/a) \\
External & one-sided vendor lane & completeness rule of the evaluation; partial measurement is reported & rule set; Intel lane: cross-check not performed \\
Conclusion & percentiles synthetic instead of measured, means distorted & nearest rank without interpolation; geometric aggregation & implemented (evaluation tree); copies in the tools: declared gap \\
Conclusion & zero instead of non-observation & field flags per counter, n/a end-to-end & implemented \\
\bottomrule
\end{tabular}
\caption{Validity of the evaluation according to the four classes of threats (construct, internal,
external, conclusion): per threat the countermeasure in the apparatus and its implementation status
--- implemented, declared and not measured (n/a), or not implemented. No row names a number without
a run.}\label{tab:eval-gueltigkeit}
\end{table}

\subsection{Summary of the Evaluation}\label{sec:eval-summary}
The evaluation records per evaluation object. Design space: The hypotheses H1--H3 carry measured
quantity and criterion, and the sensitivity analysis is fixed as a procedure over the three
measurement levels; evidence is available for the smoke series (43 permutations, seven quantities,
Appendix~\ref{app:measurements}), for the mandatory measurement series not (n/a, full build not
carried out). Apparatus: The counter-proof is fixed as a fourth dimension with test plan, validity
criterion, and decision rule, its measurement data are not available; the measurement-error
foundation --- gate fields in the fingerprint, field flags per counter, n/a rule --- is implemented.
Automation: The chain from the experiment XML to the included appendix is demonstrated in a first
end-to-end run and delivers every number with its run. The results are discussed against the
leitmotif of the thesis --- the transition from heuristic to informed cache management, that is, the
question whether telemetry-driven, automatically adapted configurations outperform statically chosen
ones; the cache-line effect from the synthesis curves is the first evaluation step after the full
run, data-dependent and therefore not anticipated here. By the standards for research
contributions~\cite{shaw2003writing} and for empirical assessments~\cite{blackburn2016truth}, the
contribution of this chapter is a validated experimental setup together with its automation, whose
result figures will be available with the full measurement series. Chapter~\ref{ch:conclusion}
answers the research questions on this basis.
